% Homework: Real-Time Operating Systems.
% Self-contained -- no Python data generation needed (the one diagram is
% drawn directly with TikZ). Compile normally for the student copy, or with
% \def\keymode{1} before \input (or -jobname + the pretex trick, same as the
% other homeworks in this repo) to produce the instructor key.
%
% By design this homework stays away from any particular RTOS's API syntax
% (no FreeRTOS, CMSIS-RTOS2, osThreadNew(), HAL_*, etc.). All pseudocode uses
% generic, vendor-neutral function names. The goal is to build conceptual
% understanding of why RTOSes exist and how they work: busy-waiting versus
% cooperative versus preemptive scheduling, threads and context switching,
% determinism, and synchronization primitives (queues, semaphores, mutexes,
% and priority inversion).
\documentclass[11pt]{article}
\usepackage[margin=0.85in]{geometry}
\usepackage{enumitem}
\usepackage{xcolor}
\usepackage{tikz}
\usetikzlibrary{positioning,arrows.meta}

\ifdefined\keymode
  \newcommand{\keyans}[1]{\par\smallskip{\color{red}\small\textbf{Key:} #1}\par}
  \newcommand{\docsuffix}{ -- Instructor Key}
\else
  \newcommand{\keyans}[1]{}
  \newcommand{\docsuffix}{}
\fi

% Ruled blank space students can use to write their answer.
\newcommand{\answerspace}[1]{\par\vspace{#1}\noindent\hrulefill\par\vspace{4pt}}

\setlength{\parindent}{0pt}
\setlist[enumerate]{leftmargin=*, itemsep=4pt, topsep=2pt}

% Priority-inversion timeline used by Problem 5: TaskL acquires a shared
% mutex, TaskH becomes ready and blocks waiting on that same mutex, and
% TaskM (which never touches the mutex at all) preempts TaskL in the
% meantime, delaying TaskH even though nothing of higher priority than
% TaskH is actually running.
\newcommand{\priorityInversionDiagram}{%
\begin{center}
\begin{tikzpicture}[x=1.55cm, y=0.95cm, >={Stealth[length=2mm]}]
  \foreach \x in {0,1,2,3,4,5}
    \draw[gray!35, dotted] (\x,0.2) -- (\x,3.9);
  \draw[->] (0,0.2) -- (5.7,0.2) node[right, font=\footnotesize] {time};
  \node[anchor=east, font=\small] at (-0.15,3) {TaskH (high)};
  \node[anchor=east, font=\small] at (-0.15,2) {TaskM (medium)};
  \node[anchor=east, font=\small] at (-0.15,1) {TaskL (low)};
  % TaskL runs and acquires the shared mutex at t0.
  \draw[line width=5pt, orange] (0,1) -- (1,1);
  % TaskH becomes ready at t1 but cannot run: it blocks waiting on the
  % mutex that TaskL is holding, all the way until t4.
  \draw[line width=5pt, red!35] (1,3) -- (4,3);
  % TaskM becomes ready at t1 as well and, having higher priority than
  % TaskL (and no interest in the mutex at all), preempts TaskL immediately.
  \draw[line width=5pt, blue] (1,2) -- (3,2);
  % TaskL only resumes once TaskM is done, then finishes its critical
  % section and releases the mutex at t4.
  \draw[line width=5pt, orange] (3,1) -- (4,1);
  % Only now can TaskH finally acquire the mutex and run.
  \draw[line width=5pt, red] (4,3) -- (5.4,3);
  \node[font=\scriptsize, align=center] at (0.9,3.5)
    {t1: TaskH ready,\\blocks on mutex};
  \node[font=\scriptsize, align=center] at (2,2.5)
    {t1--t3: TaskM preempts\\TaskL (mutex still held)};
  \node[font=\scriptsize, align=center] at (4.7,3.5)
    {t4: mutex released,\\TaskH finally runs};
\end{tikzpicture}
\end{center}}

\begin{document}

{\Large\bfseries ECGR 4101/5101 -- Advanced Embedded Systems}\hfill
Name: \underline{\hspace{5cm}}\quad Date: \underline{\hspace{2.5cm}}\\[2pt]
{\large Homework: Real-Time Operating Systems\docsuffix}

\smallskip
This homework is about the concepts behind real-time operating systems, not
the syntax of any particular one. You will \textbf{not} be asked to recall
specific API function names, structs, or configuration macros from FreeRTOS,
CMSIS-RTOS2, or any other RTOS. Any pseudocode below uses plain, made-up
function names. Answer every question in your own words, in a sentence or
two unless a question says otherwise, based on what we covered in class.

\bigskip
\textbf{Problem 1: From Busy-Waiting to a Scheduler}

\smallskip
Consider an embedded application controlling a small conveyor belt. Among
other things, it needs to pulse a stepper motor forward once every
500~ms to keep the belt moving, and it also needs to react as quickly as
possible whenever an operator presses a momentary safety-stop button wired
to a digital input pin. A first attempt at this is written as:

\begin{verbatim}
void motorStepTask(void) {
  while (getTick() - lastStepTime < 500);   // busy-wait right here
  stepMotorForward();
  lastStepTime = getTick();
}

void stopButtonTask(void) {
  if (stopButtonPressed())
    haltMotor();
}

while (1) {
  motorStepTask();
  stopButtonTask();
}
\end{verbatim}

\begin{enumerate}[label=(\alph*)]

\item Explain why a quick press-and-release of the stop button is likely to
  go completely undetected if it happens while \texttt{motorStepTask()} is
  busy-waiting, even though \texttt{stopButtonTask()} is called again and
  again inside the same \texttt{while(1)} loop.
  \answerspace{1.8cm}
  \keyans{While \texttt{motorStepTask()} is stuck in its busy-wait loop,
    control never returns to the \texttt{while(1)} loop, so
    \texttt{stopButtonTask()} cannot run at all for up to 500~ms at a time.
    If the operator presses and releases the button entirely within that
    window, the input pin has already returned to its unpressed level by
    the time \texttt{stopButtonTask()} finally checks it, so the press is
    never seen at all -- for a safety-related function like this, that is a
    serious problem, not just an inconvenience.}

\item Suppose \texttt{motorStepTask()} is rewritten so that it checks the
  elapsed time and returns immediately if 500~ms have not yet passed,
  stepping the motor and returning only once they have -- without ever
  busy-waiting inside the function. What is this general style of
  scheduling called, and what does it mean, conceptually, for a task to
  ``cooperate'' with the rest of the system?
  \answerspace{1.6cm}
  \keyans{This is cooperative scheduling. ``Cooperating'' means each task is
    written so that it voluntarily and quickly gives up the CPU (by
    returning or otherwise yielding) rather than holding onto it, trusting
    that the other tasks will do the same.}

\item This rewritten \texttt{motorStepTask()} still needs a variable like
  \texttt{lastStepTime} to remember the tick count from the last time it
  ran. Explain why this value cannot simply be a local variable declared
  inside \texttt{motorStepTask()}, and what this implies about what a
  ``thread'' running under an RTOS needs to have of its own that a plain
  function call does not.
  \answerspace{1.8cm}
  \keyans{A local variable's storage (the function's stack frame) is
    destroyed the moment the function returns, so a local
    \texttt{lastStepTime} would be reset to garbage every single call
    instead of remembering the previous step time. A thread needs its own
    persistent storage -- in particular, its own stack and saved execution
    context -- that survives across the many separate ``turns'' it is
    given, which is exactly what an RTOS sets up for each thread it
    creates.}

\item Now suppose the stop button is instead wired to trigger a hardware
  interrupt the instant it is pressed, and the interrupt handler itself
  halts the motor (rather than the main loop periodically asking ``has the
  button been pressed?''). What general term describes this style of
  handling the stop-button task compared to the cooperative version in part
  (b), and why does it no longer matter how long \texttt{motorStepTask()}
  takes between checks, or how brief the button press is?
  \answerspace{1.8cm}
  \keyans{This is preemptive scheduling/servicing: the button press itself
    forcibly interrupts whatever is currently running, rather than waiting
    for that code to voluntarily check whether the button is pressed.
    Because the hardware guarantees the handler runs immediately on the
    press, it no longer matters how long \texttt{motorStepTask()} takes
    between checks, nor whether the press-and-release happens faster than
    the main loop could have polled for it.}

\item As an application grows to include many tasks with very different
  timing requirements (some periodic, some triggered by external events,
  some with strict deadlines), explain why hand-coding all of their
  interactions inside one big \texttt{while(1)} loop -- even a carefully
  cooperative one -- eventually becomes impractical. What part of an RTOS
  exists specifically to take over this responsibility?
  \answerspace{1.8cm}
  \keyans{With only a handful of tasks it is feasible for a programmer to
    reason about every possible ordering and timing interaction by hand,
    but as the number of tasks grows, the number of possible interleavings
    and worst-case delays grows far faster, making it increasingly
    difficult (and eventually impossible in practice) to guarantee by
    inspection that no task will miss a deadline or block another for too
    long. The scheduler is the part of an RTOS that takes over the job of
    deciding, at every moment, which ready task should run next.}

\end{enumerate}

\newpage
\textbf{Problem 2: Threads, Context Switches, and Scheduling Policies}

\begin{enumerate}[label=(\alph*)]

\item In your own words, define what a ``thread'' is in the context of an
  RTOS running on a single-core microcontroller.
  \answerspace{1.4cm}
  \keyans{A thread is an independent stream of execution -- conceptually,
    just a function with its own associated state -- that the scheduler can
    run, pause, and resume independently of any other thread, giving the
    illusion that it executes concurrently with the rest of the
    application.}

\item A single CPU core can only execute instructions from one thread at
  any given instant, yet several threads appear to make progress at the
  same time. Explain the mechanism that creates this illusion, and name the
  specific action the OS performs to stop running one thread and begin
  running another.
  \answerspace{1.6cm}
  \keyans{The illusion is created by rapidly switching which thread's
    instructions the CPU is executing, giving each one a short slice of
    time so that, to an outside observer, all of them seem to progress at
    once even though only one is ever truly running at a time. The specific
    action that swaps the CPU from one thread's state to another's is
    called a context switch.}

\item During a context switch, the RTOS must save and later restore each
  thread's execution state in memory dedicated to that specific thread.
  List at least three distinct pieces of information that must be saved,
  and explain why each thread needs its own separate block of memory to
  hold this information rather than sharing one block among all threads.
  \answerspace{1.8cm}
  \keyans{At minimum, the CPU registers (including the program counter and
    general-purpose registers), the thread's own stack pointer, and the
    contents of that thread's own stack (its local variables and call
    history) all need to be preserved. Each thread needs its own block of
    memory because its registers and stack contents will, in general, be
    completely different from every other thread's at the moment it gets
    paused -- overwriting them with another thread's values (or sharing one
    stack among threads) would destroy the paused thread's ability to
    resume correctly later.}

\item Describe, in your own words, the difference between round-robin
  scheduling, cooperative scheduling, and preemptive priority-based
  scheduling. For each one, state specifically what determines when a
  currently running thread stops running and a different thread starts.
  \answerspace{2.2cm}
  \keyans{Round-robin: every ready thread is given an equal, fixed slice of
    time in turn, and a thread stops running simply because its time slice
    has expired, regardless of priority. Cooperative: a thread keeps
    running until it voluntarily gives up the CPU (by blocking, finishing,
    or explicitly yielding) -- the scheduler never forcibly interrupts it.
    Preemptive priority-based: a running thread is forcibly stopped the
    instant a higher-priority thread becomes ready to run, regardless of
    whether the running thread was willing to give up the CPU or not.}

\item Suppose a low-priority thread enters an infinite loop that never
  calls any blocking function and never returns. Under a preemptive
  scheduler, explain what hardware mechanism still allows a
  higher-priority thread to get CPU time, and what would happen to the
  rest of the system if that hardware mechanism did not exist.
  \answerspace{1.8cm}
  \keyans{A periodic hardware timer interrupt (separate from anything the
    misbehaving thread does) forces control away from the currently
    running thread at a fixed interval, giving the scheduler a chance to
    run and switch to a higher-priority ready thread. Without this
    mechanism, a thread that never voluntarily returns control would run
    forever, permanently starving every other thread in the system -- no
    higher-priority thread, no matter how urgent, would ever get to run
    again.}

\item Every RTOS includes a special, built-in ``idle'' thread that is
  always ready to run and is assigned the lowest possible priority.
  Explain why an OS needs a thread like this instead of simply leaving the
  CPU with nothing scheduled whenever every other thread is blocked or
  suspended.
  \answerspace{1.6cm}
  \keyans{The scheduler always needs some thread to hand the CPU to; if
    every application thread happens to be blocked or suspended at the same
    moment, there still has to be something for the processor to execute.
    The idle thread exists purely to fill this need (and is commonly also
    used to perform low-priority housekeeping, such as freeing resources
    of deleted threads) whenever no other thread is ready.}

\end{enumerate}

\newpage
\textbf{Problem 3: Real-Time Behavior and Determinism}

\begin{enumerate}[label=(\alph*)]

\item Explain what specifically distinguishes a ``real-time'' operating
  system from a general-purpose operating system, in terms of what
  guarantee it must provide and over what kind of time scale that
  guarantee is expected to hold.
  \answerspace{1.6cm}
  \keyans{A real-time OS must guarantee that certain operations complete
    within a specified time constraint, called a deadline, which is
    typically on the order of milliseconds or even microseconds. A
    general-purpose OS has no such guarantee -- it may respond quickly most
    of the time, but nothing in its design promises an upper bound on how
    long a response might take.}

\item Desktop operating systems rely heavily on virtual memory, paging, and
  swapping data out to disk. Explain, in your own words, why page faults
  and disk-based swapping are fundamentally difficult to reconcile with
  meeting a hard real-time deadline, even though this behavior works
  perfectly well for everyday desktop applications.
  \answerspace{1.8cm}
  \keyans{Servicing a page fault means pausing the running program,
    locating the needed data (possibly on a disk), and loading it back into
    RAM before execution can continue -- and the amount of time this takes
    is not fixed; it depends on where the data physically lives and on
    whatever else the storage device happens to be doing at that moment.
    A system that cannot guarantee how long this pause will last cannot
    guarantee that any deadline will be met, which is why real-time systems
    generally avoid paging and swapping even though ordinary desktop use
    rarely notices the resulting delay.}

\item Explain why a thread's priority (used by the scheduler to decide
  which ready thread runs next) is a fundamentally different concept from
  the priority assigned to a hardware interrupt, and give an example of a
  situation where this distinction actually matters -- that is, a moment
  when a high-priority thread is ready to run but is not actually the thing
  executing on the CPU.
  \answerspace{1.8cm}
  \keyans{Thread priority only governs competition among threads for the
    CPU once an interrupt has finished and the scheduler is choosing which
    thread to resume; it says nothing about what happens while an interrupt
    service routine itself is executing. For example, while any interrupt
    service routine is running (even one unrelated to a given high-priority
    thread), that thread is formally ``ready'' and may even be the highest
    priority ready thread in the whole system, yet it is still not the
    thing actually executing on the CPU until the ISR finishes and the
    scheduler runs.}

\end{enumerate}

\newpage
\textbf{Problem 4: Synchronization with Queues and Semaphores}

\smallskip
Suppose one thread continually produces data (for example, samples read
from a sensor) and a second thread continually consumes that data (for
example, computing a running average) by sharing a common buffer between
them.

\begin{enumerate}[label=(\alph*)]

\item Describe, in your own words, the two mistakes that must both be
  avoided for this producer/consumer arrangement to work correctly: one
  mistake the producer could make, and one mistake the consumer could
  make.
  \answerspace{1.6cm}
  \keyans{The producer must not keep writing new data into the buffer once
    it is completely full (overwriting data the consumer has not yet
    read). The consumer must not try to read data out of the buffer once
    it is completely empty (reading stale or meaningless data that was
    never actually produced).}

\item Explain how a message queue, as provided by an RTOS, prevents both
  of those mistakes automatically, including what specifically happens to
  a thread that tries to add an item to a full queue, or to a thread that
  tries to remove an item from an empty queue.
  \answerspace{1.8cm}
  \keyans{A message queue keeps track of how many items it currently holds
    and enforces its fixed capacity. A thread that tries to add an item
    to a full queue is placed in the blocked state until at least one item
    is removed by some other thread, and likewise a thread that tries to
    remove an item from an empty queue is placed in the blocked state until
    some other thread adds a new item -- in both cases the OS kernel
    automatically pauses the offending thread rather than letting it
    proceed incorrectly.}

\item A semaphore can conceptually be thought of as nothing more than a
  variable (a boolean flag, or a counter) that threads check before
  proceeding. Explain why a plain, unprotected global variable cannot
  safely serve as a semaphore in a multithreaded system, and what role the
  OS kernel has to play in order to make this safe.
  \answerspace{1.8cm}
  \keyans{There is no guarantee that checking a variable's value and then
    updating it happens as a single, uninterruptible step; a thread could
    be preempted in between checking and updating the variable, allowing
    another thread to also ``succeed'' at the same check before the first
    thread finishes updating it, letting two threads believe they both
    safely acquired the same resource. The OS kernel must perform the
    check-and-update as an atomic operation (for example, by disabling
    interrupts or otherwise preventing a context switch during that brief
    step) so that only one thread can ever succeed at acquiring the
    semaphore at a time.}

\item Explain the conceptual difference between a binary semaphore and a
  counting semaphore. Then describe one resource-sharing scenario of your
  own (not one used in this course) where a counting semaphore with an
  initial count greater than one would be the appropriate choice.
  \answerspace{1.8cm}
  \keyans{A binary semaphore can only be in one of two states (taken or
    available), so it is suited to situations where exactly one thread may
    hold the resource at a time. A counting semaphore starts with some
    initial count greater than one and is decremented every time it is
    acquired and incremented every time it is released, so up to that many
    threads may hold instances of the resource simultaneously, and a
    thread only blocks once the count reaches zero. Any reasonable example
    is acceptable, e.g., a pool of several identical hardware DMA channels,
    a fixed number of reusable network connection buffers, or a small pool
    of identical sensor-reading structures, where a counting semaphore
    initialized to the number of available resources in the pool lets that
    many threads proceed at once.}

\end{enumerate}

\newpage
\textbf{Problem 5: Mutexes and Priority Inversion}

\smallskip
The timeline below shows three threads of different priority:
\texttt{TaskH} (high priority), \texttt{TaskM} (medium priority), and
\texttt{TaskL} (low priority). \texttt{TaskL} and \texttt{TaskH} both need
to access the same shared resource and therefore both acquire the same
mutex before using it; \texttt{TaskM} never touches that resource or that
mutex at all.

\priorityInversionDiagram

\begin{enumerate}[label=(\alph*)]

\item Using the diagram, describe in your own words the sequence of events
  that causes the high-priority \texttt{TaskH} to sit blocked for so long,
  even though no thread with higher priority than \texttt{TaskH} is running
  anywhere in the system during that time. This general problem is known
  as priority inversion.
  \answerspace{2cm}
  \keyans{\texttt{TaskL} acquires the mutex first. Shortly after,
    \texttt{TaskH} becomes ready and tries to acquire that same mutex, but
    since \texttt{TaskL} is still holding it, \texttt{TaskH} must block and
    wait. Meanwhile, \texttt{TaskM} becomes ready at around the same time
    and, having higher priority than \texttt{TaskL} (and no interest in the
    mutex at all), immediately preempts \texttt{TaskL} and runs to
    completion. Only once \texttt{TaskM} finishes does \texttt{TaskL}
    resume, finish its critical section, and finally release the mutex --
    at which point \texttt{TaskH} can finally acquire it and run. The net
    effect is that the medium-priority \texttt{TaskM}, despite never
    touching the shared resource, has effectively delayed the
    highest-priority thread in the system.}

\item Explain the conceptual difference between a mutex and a semaphore in
  terms of ownership: who is expected to acquire and release each one.
  \answerspace{1.6cm}
  \keyans{A mutex is meant to always be acquired and later released by that
    same thread, since it exists purely to protect a shared resource during
    the period that one specific thread is using it. A semaphore has no
    such ownership expectation -- it is entirely normal, and often by
    design, for one thread (or even an interrupt handler) to release a
    semaphore that a completely different thread later acquires, since
    semaphores are also used to signal that some event has occurred rather
    than only to protect a resource.}

\item Describe the general strategy (sometimes called priority
  inheritance) that many RTOSes use to minimize priority inversion.
  Explain, specifically, how temporarily changing the priority of one of
  the three threads in the diagram above would have prevented
  \texttt{TaskM} from delaying \texttt{TaskL} (and therefore
  \texttt{TaskH}) for as long as it did.
  \answerspace{2cm}
  \keyans{When a higher-priority thread is found to be blocked waiting on a
    mutex that a lower-priority thread currently holds, the OS temporarily
    raises the priority of the mutex-holding thread to match the priority
    of the thread waiting on it. In the diagram, this means that as soon
    as \texttt{TaskH} blocks waiting on the mutex held by \texttt{TaskL},
    \texttt{TaskL}'s priority would be temporarily raised to \texttt{TaskH}'s
    level. \texttt{TaskM} would then no longer be able to preempt
    \texttt{TaskL}, since \texttt{TaskL} now effectively outranks it, so
    \texttt{TaskL} would finish its critical section and release the mutex
    much sooner, letting \texttt{TaskH} run far earlier -- and only after
    the mutex is released would \texttt{TaskL}'s priority drop back down,
    finally allowing \texttt{TaskM} to run.}

\item Suppose priority inheritance were not implemented at all in this
  example. Explain, in the worst case, how additional medium-priority
  threads could prevent the high-priority \texttt{TaskH} from running for
  an arbitrarily long, essentially unbounded amount of time -- a phenomenon
  sometimes called unbounded priority inversion.
  \answerspace{1.8cm}
  \keyans{Without priority inheritance, \texttt{TaskL}'s priority never
    changes while it holds the mutex, so it remains just as vulnerable to
    being preempted by any medium-priority thread as it was before
    acquiring the mutex. If new medium-priority threads keep becoming ready
    one after another, each one will preempt \texttt{TaskL} in turn before
    it can finish its critical section and release the mutex, and
    \texttt{TaskH} remains blocked the entire time. Since there is no limit
    on how many such medium-priority threads could keep arriving, there is
    also no bound on how long \texttt{TaskH} -- the highest-priority thread
    in the system -- could be forced to wait.}

\end{enumerate}

\end{document}
